% Persian manual of the pcalendar package. Build: make doc   (XeLaTeX, font XB Niloofar)
\documentclass[11pt]{article}

\usepackage[a4paper,margin=2.5cm]{geometry}
\usepackage{pcalendar}
\usepackage{pgffor}
\usepackage[colorlinks,linkcolor=blue!50!black,urlcolor=blue!50!black]{hyperref}
\usepackage{xepersian}

\settextfont{XB Niloofar}
\setdigitfont{Parisa}
\setlatintextfont[Extension=.otf,BoldFont=lmroman10-bold,ItalicFont=lmroman10-italic,BoldItalicFont=lmroman10-bolditalic]{lmroman10-regular}

% --- helpers --------------------------------------------------------------
\newcommand\tx[1]{\mbox{\lr{\texttt{#1}}}}                  % code inside Persian text
\newcommand\cs[1]{\mbox{\lr{\texttt{\textbackslash #1}}}}   % a control sequence
\newcommand\pkg[1]{\mbox{\lr{\textsf{#1}}}}
\newcommand\iso[1]{\mbox{\lr{#1}}}                          % numeric dates must stay left to right
\newcommand\gap{\par\smallskip}
% \macro|syntax|: heading line for a macro description
\NewDocumentCommand\macro{v}{\par\bigskip\noindent\lr{\texttt{\bfseries #1}}\par\nopagebreak\smallskip}
% \ex|code|{result}: a call, an arrow and its typeset result (the call is typed twice on purpose)
\NewDocumentCommand\ex{v m}{\par\noindent\hspace*{1em}\lr{\texttt{#1}}\quad$\leftarrow$\quad#2}
\emergencystretch=3em
\newcount\jd \newcount\wd
\makeatletter\renewcommand*\l@section{\@dottedtocline{1}{0em}{2.3em}}\makeatother
% table cells for the name lists: \mc{n} month n, \yc{y} year y; \MROW/\YROW{first} fill a row of four
\newcommand\mc[1]{\edef\mm{\number\numexpr#1\relax}\tx{\mm}\quad\pcalendarmonthname{\mm}}
\newcommand\yc[1]{\edef\yy{\number\numexpr#1\relax}\tx{\yy}\quad\pcalendaryearname{\yy}}
\newcommand\MROW[1]{\mc{#1}&\mc{#1+3}&\mc{#1+6}&\mc{#1+9}}
\newcommand\YROW[1]{\yc{#1}&\yc{#1+3}&\yc{#1+6}&\yc{#1+9}}

\title{بستهٔ \pkg{pcalendar}}
\author{جواد رضویان\\\tx{javadr@gmail.com}}
\date{نسخهٔ \lr{\pcalendarversion}\quad(\lr{\pcalendardate})}

\begin{document}

\maketitle

\begin{abstract}\noindent
بستهٔ \pkg{pcalendar} تاریخ را بین تقویم جلالی (هجری شمسی) و میلادی تبدیل می‌کند و آن را با نام ماه‌های
فارسی چاپ می‌کند. با کمک شمارهٔ روز ژولین می‌توانید به تاریخ‌ها چند روز اضافه یا از آن‌ها کم کنید. نام
ماه‌ها، نام روزهای هفته و نام سال‌های دوازده‌حیوانی هم در بسته آماده است. \pkg{pcalendar} بر پایهٔ
\pkg{pgfcalendar} ساخته شده و برای کار با \pkg{xepersian} در \lr{XeLaTeX} نوشته شده است.
\end{abstract}

\tableofcontents

% ---------------------------------------------------------------------------
\section{معرفی}\label{sec:what}

\pkg{pcalendar} مبدل تاریخ است و چند جدول نام هم دارد. با آن می‌توانید:
\begin{itemize}
  \item تاریخ میلادی را به جلالی و تاریخ جلالی را به میلادی تبدیل کنید، یا به‌صورت متن (مثل «۱ فروردین
    ۱۳۹۴») یا به‌صورت عدد (مثل ۰۱/۰۱/۱۳۹۴)؛
  \item شمارهٔ روز ژولین یک تاریخ جلالی را بگیرید، چند روز به آن اضافه یا از آن کم کنید و دوباره به تاریخ
    برگردانید؛
  \item نام ماه‌های جلالی و میلادی و نام روزهای هفته را چاپ کنید؛
  \item ببینید هر سال به کدام حیوان از دورهٔ دوازده‌ساله تعلق دارد.
\end{itemize}
در مقابل، این بسته جدول ماه یا تقویم رومیزی نمی‌سازد، از مناسبت‌ها و تعطیلات رسمی خبر ندارد و فرمان
جداگانه‌ای هم برای روز هفتهٔ یک تاریخ ندارد. با این حال، در بخش~\ref{sec:names} می‌بینید که با چند خط
می‌توان روز هفته را به دست آورد.

\paragraph{پیش‌نیازها.} \lr{XeLaTeX}، بستهٔ \pkg{pgf} (که \pkg{pgfcalendar} بخشی از آن است) و
\pkg{xepersian} همراه با یک قلم فارسی لازم است. نام‌های فارسی به‌صورت متن \lr{UTF-8} در خود فایل بسته نوشته
شده‌اند.

% ---------------------------------------------------------------------------
\section{نصب}\label{sec:install}

اگر توزیع \TeX{} شما این بسته را داشته باشد، لازم نیست کاری بکنید. برای نصب دستی، فایل \tx{pcalendar.sty} را
در یکی از پوشه‌هایی که \TeX{} فایل‌ها را در آن‌ها جست‌وجو می‌کند قرار دهید، مثلاً
\tx{texmf/tex/latex/pcalendar/}، و با \tx{mktexlsr} فهرست فایل‌ها را به‌روز کنید. در مخزن پروژه، دستور
\tx{make dev} همین کار را انجام می‌دهد. اگر نمی‌خواهید چیزی نصب کنید، همین‌قدر کافی است که
\tx{pcalendar.sty} را کنار سند خودتان بگذارید.

بسته را در مقدمهٔ سند فراخوانی کنید. در این راهنما، \pkg{pcalendar} پیش از \pkg{xepersian} بارگذاری می‌شود:
\begin{latin}\begin{verbatim}
\usepackage{pcalendar}
\usepackage{xepersian}
\settextfont{XB Niloofar}
\end{verbatim}\end{latin}
بسته گزینه‌ای ندارد.

% ---------------------------------------------------------------------------
\section{شروع سریع}\label{sec:quick}

\begin{latin}\begin{verbatim}
\pdate{2015-03-21}          % Gregorian in, Jalali text out
\gdate{1394-01-01}          % Jalali in, Gregorian text out
\pisodate{2015-03-21}       % Gregorian in, Jalali as DD/MM/YYYY
\gisodate{1394-01-01}       % Jalali in, Gregorian as DD/MM/YYYY
\end{verbatim}\end{latin}
خروجی این چهار فرمان به این صورت است:
\ex|\pdate{2015-03-21}|{\pdate{2015-03-21}}
\ex|\gdate{1394-01-01}|{\gdate{1394-01-01}}
\ex|\pisodate{2015-03-21}|{\iso{\pisodate{2015-03-21}}}
\ex|\gisodate{1394-01-01}|{\iso{\gisodate{1394-01-01}}}
\gap
تاریخ امروز را هم می‌توان به تقویم جلالی چاپ کرد. اگر
\cs{pdate}\tx{\{\textbackslash year-\textbackslash month-\textbackslash day\}} را بنویسید، امروز چنین
چاپ می‌شود: \pdate{\year-\month-\day}.

% ---------------------------------------------------------------------------
\section{تبدیل تاریخ}\label{sec:date}

هر چهار فرمان یک ورودی دارند: تاریخی به شکل \tx{Y-M-D}، یعنی سال، شمارهٔ ماه و شمارهٔ روز که با خط تیره از
هم جدا شده‌اند. صفر ابتدای ماه و روز اختیاری است؛ یعنی \tx{2015-3-21} با \tx{2015-03-21} یکی است.

\macro|\pdate{Y-M-D}|
ورودی تاریخ میلادی است و خروجی تاریخ جلالی به شکل «روز، نام ماه، سال»؛ مثلاً \pdate{2024-03-20}.

\macro|\gdate{Y-M-D}|
ورودی تاریخ جلالی است و خروجی تاریخ میلادی به همان شکل «روز، نام ماه، سال». نام ماه میلادی هم فارسی
نوشته می‌شود؛ مثلاً \gdate{1403-01-01}.

\macro|\pisodate{Y-M-D}|
مانند \cs{pdate} از تاریخ میلادی تاریخ جلالی می‌سازد، اما خروجی‌اش عددی است: روز و ماه دورقمی‌اند و با
اسلش از هم جدا می‌شوند؛ مثلاً \iso{\pisodate{2024-03-20}}. ترتیب خروجی روز/ماه/سال است، نه سال-ماه-روز
مانند استاندارد \lr{ISO}.

\macro|\gisodate{Y-M-D}|
مانند \cs{gdate} از تاریخ جلالی تاریخ میلادی می‌سازد، اما عددی و با همان ترتیب روز/ماه/سال؛ مثلاً
\iso{\gisodate{1403-01-01}}.

\paragraph{سال کبیسه.} اسفند در سال کبیسه ۳۰ روز دارد:
\ex|\pdate{2025-03-20}|{\pdate{2025-03-20}}
\ex|\gdate{1403-12-30}|{\gdate{1403-12-30}}
\ex|\gdate{1404-12-30}|{\gdate{1404-12-30}}
تاریخ آخر نادرست است، چون ۱۴۰۴ کبیسه نیست؛ با این حال بسته درست‌بودن ورودی را بررسی نمی‌کند
(بخش~\ref{sec:limits}).

\paragraph{نکتهٔ حروف‌چینی.} در متن فارسی ممکن است ترتیب عددها به‌هم بریزد (به‌خاطر قاعدهٔ متن دوسویه).
برای اینکه خروجی \cs{pisodate} و \cs{gisodate} همیشه به ترتیب روز/ماه/سال بماند، آن را درون \cs{lr}
بگذارید؛ همان کاری که در همین راهنما شده است.

% ---------------------------------------------------------------------------
\section{شمارهٔ روز و محاسبه روی تاریخ}\label{sec:julian}

فرمان‌های بالا فقط تبدیل می‌کنند و نمی‌توان با آن‌ها به تاریخ روز اضافه کرد. برای جمع و تفریق روز، بسته از
شمارهٔ روز ژولین استفاده می‌کند، همان شماره‌ای که \pkg{pgfcalendar} تعریف کرده است. روزها از یک مبدأ ثابت
شمرده می‌شوند؛ بنابراین تفاضل دو شماره دقیقاً برابر با تعداد روزهای بین دو تاریخ است.

\macro|\pcalendardatetojulian{Y-M-D}{\count}|
شمارهٔ روز تاریخ \emph{جلالی} \tx{Y-M-D} را در شمارندهٔ \cs{count} می‌گذارد. اگر \tx{+N} را به انتهای تاریخ
اضافه کنید، \tx{N} روز هم به آن افزوده می‌شود؛ مثلاً \tx{1394-01-01+10} یعنی ده روز بعد از اول فروردین ۱۳۹۴.

\macro|\pcalendarjuliantodate{\count}{\y}{\m}{\d}|
شمارهٔ روزِ موجود در \cs{count} را به تاریخ \emph{جلالی} برمی‌گرداند و سال، ماه و روز را در فرمان‌های \cs{y}،
\cs{m} و \cs{d} قرار می‌دهد (به‌جای این سه می‌توانید نام‌های دیگری هم بدهید).

\begin{latin}\begin{verbatim}
\newcount\jd
\pcalendardatetojulian{1394-01-01+10}\jd
\pcalendarjuliantodate\jd\y\m\d
\number\y-\number\m-\number\d
\end{verbatim}\end{latin}
\pcalendardatetojulian{1394-01-01+10}\jd%
\pcalendarjuliantodate\jd\y\m\d%
خروجی این کد \iso{\number\y-\number\m-\number\d} است.

\paragraph{\cs{y} و \cs{m} و \cs{d} را بلافاصله بخوانید.} بعد از \cs{pcalendarjuliantodate}، این سه فرمان
رونوشتی از مقدار نیستند؛ فقط نام دیگری برای شمارنده‌های نتیجهٔ بسته‌اند. بنابراین هر تبدیل بعدی
(\cs{pdate}، \cs{gisodate} و مانند آن‌ها) مقدارشان را عوض می‌کند. اگر بعداً به آن‌ها نیاز دارید، همان
لحظه در شمارنده‌های خودتان ذخیره‌شان کنید:
\tx{\textbackslash newcount\textbackslash myyear\ \textbackslash myyear=\textbackslash y}.

\paragraph{عبور از پایان سال.} این محاسبه سال کبیسه را هم درنظر می‌گیرد:
\pcalendardatetojulian{1403-12-29+1}\jd\pcalendarjuliantodate\jd\y\m\d%
یک روز بعد از ۱۴۰۳-۱۲-۲۹ می‌شود \iso{\number\y-\number\m-\number\d}، چون ۱۴۰۳ کبیسه است؛
\pcalendardatetojulian{1403-12-29+2}\jd\pcalendarjuliantodate\jd\y\m\d%
و دو روز بعد از آن می‌شود \iso{\number\y-\number\m-\number\d}.

\paragraph{شمارنده‌های موقت.} \cs{pcalendardatetojulian} از \cs{count1}، \cs{count2}، \cs{count3} و
\cs{count8} به‌عنوان فضای کار استفاده می‌کند. مقداری در این شمارنده‌ها نگه ندارید.

% ---------------------------------------------------------------------------
\section{نام‌ها}\label{sec:names}

\macro|\pcalendarmonthname{n}|
نام ماه جلالی را چاپ می‌کند. \tx{n} عددی از ۱ تا ۱۲ است و برای عددهای بیرون از این بازه چیزی چاپ نمی‌شود.
\ex|\pcalendarmonthname{1}|{\pcalendarmonthname{1}}
\gap
نام دوازده ماه، به ترتیب:\gap
\begin{tabular}{llll}
\MROW{1}\\
\MROW{2}\\
\MROW{3}
\end{tabular}

\macro|\pcalendargregorianmonthname{n}|
نام ماه میلادی را به فارسی چاپ می‌کند (مثلاً «مارس»). \tx{n} عددی از ۱ تا ۱۲ است.
\ex|\pcalendargregorianmonthname{3}|{\pcalendargregorianmonthname{3}}

\macro|\pcalendarweekdayname{n}|
نام روز هفته را چاپ می‌کند. \tx{n} عددی از ۰ تا ۶ است: ۰ شنبه (اولین روز هفته در ایران) و ۶ جمعه است.
\gap\leavevmode
\foreach \w in {0,...,6}{\hbox{\tx{\w}~\pcalendarweekdayname{\w}}\quad}

\macro|\pcalendarweekdayshortname{n}|
مخفف تک‌حرفی نام روز هفته را چاپ می‌کند؛ شماره‌گذاری مانند فرمان قبلی است.
\gap\leavevmode
\foreach \w in {0,...,6}{\hbox{\tx{\w}~\pcalendarweekdayshortname{\w}}\quad}

\macro|\pcalendaryearname{year}|
نام یک سال جلالی را در دورهٔ دوازده‌حیوانی چاپ می‌کند. \tx{year} باید بین ۰ تا ۲۰۹۹ باشد و برای مقدارهای
دیگر چیزی چاپ نمی‌شود.
\ex|\pcalendaryearname{1403}|{\pcalendaryearname{1403}}
\gap
دوازده حیوان، به ترتیب و از سال ۱۳۸۰:\gap
\begin{tabular}{llll}
\YROW{1380}\\
\YROW{1381}\\
\YROW{1382}
\end{tabular}

\paragraph{روز هفتهٔ یک تاریخ.} فرمان آماده‌ای برای این کار وجود ندارد، اما از روی شمارهٔ روز می‌توان آن را
به دست آورد. \pkg{pgfcalendar} دوشنبه را ۰ می‌گیرد و \cs{pcalendarweekdayname} شنبه را؛ پس نتیجه را دو
واحد جابه‌جا می‌کنیم:
\begin{latin}\begin{verbatim}
\newcount\jd \newcount\wd
\pcalendardatetojulian{1403-01-01}\jd
\pgfcalendarjuliantoweekday{\jd}{\wd}
\advance\wd by 2 \ifnum\wd>6 \advance\wd by -7 \fi
\pcalendarweekdayname{\wd}
\end{verbatim}\end{latin}
\pcalendardatetojulian{1403-01-01}\jd%
\pgfcalendarjuliantoweekday{\jd}{\wd}%
\advance\wd by 2 \ifnum\wd>6 \advance\wd by -7 \fi%
نتیجه «\pcalendarweekdayname{\wd}» است؛ یعنی اول فروردین ۱۴۰۳ چهارشنبه بوده است.

% ---------------------------------------------------------------------------
\section{نمونه‌های کاربردی}\label{sec:recipes}

\paragraph{فاصلهٔ دو تاریخ جلالی به روز.}\leavevmode\par
\begin{latin}\begin{verbatim}
\newcount\a \newcount\b
\pcalendardatetojulian{1403-01-01}\a
\pcalendardatetojulian{1404-01-01}\b
\advance\b by -\a
\number\b
\end{verbatim}\end{latin}
\pcalendardatetojulian{1403-01-01}\jd\newcount\ja\ja=\jd
\pcalendardatetojulian{1404-01-01}\jd\advance\jd by -\ja
نتیجه \number\jd{} روز است (۱۴۰۳ کبیسه است).

\paragraph{چاپ چند روز پیاپی.} برای چاپ پنج روز از اول فروردین ۱۴۰۳ می‌توان از حلقه استفاده کرد. هر تاریخ
باید همان‌جا تبدیل و چاپ شود، چون نتیجه با تبدیل بعدی عوض می‌شود:
\begin{latin}\begin{verbatim}
\newcount\jd
\foreach \x in {0,...,4}{%
  \pcalendardatetojulian{1403-01-01+\x}\jd
  \pcalendarjuliantodate\jd\y\m\d
  \number\y-\number\m-\number\d\par
}
\end{verbatim}\end{latin}
\foreach \x in {0,...,4}{%
  \pcalendardatetojulian{1403-01-01+\x}\jd
  \pcalendarjuliantodate\jd\y\m\d
  \par\noindent\hspace*{1em}\iso{\number\y-\number\m-\number\d}%
}
\par\medskip
استفاده از \tx{\textbackslash x} درون \cs{pcalendardatetojulian} مشکلی ندارد، چون این فرمان ورودی‌اش را باز
می‌کند. فایل نمونهٔ \tx{phd-study.tex} در مخزن پروژه از همین روش برای شماره‌گذاری ۱۴۵ روز یک برنامهٔ مطالعه
استفاده کرده است.

\paragraph{شمارنده‌ها را بیرون از حلقه تعریف کنید.} \cs{newcount} درون حلقه در هر دور یک شمارندهٔ تازه
برمی‌دارد و سرانجام کار با خطای «\lr{No room for a new \textbackslash count}» متوقف می‌شود.

% ---------------------------------------------------------------------------
\section{سازوکار داخلی}\label{sec:how}

تبدیل‌ها با حساب ساده روی شمارنده‌های \TeX{} انجام می‌شوند و نه به برنامهٔ بیرونی نیاز دارند، نه به \lr{Lua}.
کار اصلی را دو فرمان داخلی انجام می‌دهند. هر فرمان عمومی یکی از این دو را صدا می‌زند و سپس نتیجه را چاپ یا
کپی می‌کند.

\begin{description}
  \item[\tx{\textbackslash pcal@gregoriantojalali\{Y-M-D\}}] تبدیل میلادی به جلالی. ابتدا کبیسه‌بودن سال
    میلادی را تعیین و جدول طول ماه‌های هر دو تقویم را آماده می‌کند. سپس تعداد روزهای گذشته از ابتدای سال
    میلادی تا تاریخ ورودی را می‌شمارد و آن‌قدر در ماه‌های جلالی جلو می‌رود که این روزها تمام شوند. سال‌های
    کبیسهٔ جلالی با قاعدهٔ رایج \emph{دورهٔ ۳۳ ساله} مشخص می‌شوند: در هر ۳۳ سال، هشت سال کبیسه است.
  \item[\tx{\textbackslash pcal@jalalitogregorian\{Y-M-D\}}] تبدیل جلالی به میلادی. ابتدا از روی سال، تعداد
    روزهای گذشته از یک مبدأ ثابت را حساب می‌کند (۳۶۵ روز برای هر سال، به‌اضافهٔ ۸ روز کبیسه در هر ۳۳ سال و
    روزهای کبیسهٔ دورهٔ ناقص). بعد روزهای ماه‌های قبلی را به آن اضافه می‌کند و مجموع را به بلوک‌های ۴۰۰ ساله،
    ۱۰۰ ساله، ۴ ساله و ۱ ساله میلادی تقسیم می‌کند.
\end{description}

\paragraph{شمارنده‌های نتیجه.} نتیجهٔ جلالی در \cs{pcal@y}، \cs{pcal@m} و \cs{pcal@d} و نتیجهٔ میلادی در
\cs{pcal@gy}، \cs{pcal@gm} و \cs{pcal@gd} ذخیره می‌شود. پرچم‌های کبیسه \cs{ifpcal@leap} (میلادی) و
\cs{ifpcal@kabiseh} (جلالی) هستند. طول ماه‌ها در شمارنده‌هایی نگهداری می‌شود که نامشان به عدد رومی ختم
می‌شود (\cs{pcal@latini} تا \cs{pcal@latinxii} و \cs{pcal@persiani} تا \cs{pcal@persianxiii}). دی دو بار
می‌آید: خانهٔ اول بخشی از دی است که در سال میلادی باقی مانده و خانهٔ آخر کل ماه را نشان می‌دهد؛ برای همین
سیزده خانه لازم است.

\paragraph{پل شمارهٔ روز ژولین.} \cs{pcalendardatetojulian} ابتدا تاریخ جلالی را با فرمان جلالی-به-میلادیِ
بالا به تاریخ میلادی برمی‌گرداند و آن را به \cs{pgfcalendardatetojulian} می‌دهد. \cs{pcalendarjuliantodate}
عکس این مسیر را می‌رود: \cs{pgfcalendarjuliantodate} تاریخ میلادی می‌سازد و فرمان میلادی-به-جلالی آن را
به تاریخ جلالی تبدیل می‌کند.

\paragraph{نام سال.} حیوان‌های سال هر ۱۲ سال تکرار می‌شوند. علاوه بر این، در هر سده یک جابه‌جایی سه‌حالتی
اعمال می‌شود تا سال ۰ دوره با حیوانی غیر از حیوان سال ۱۰۰ شروع شود. به همین دلیل این فرمان فقط تا ۲۰۹۹ کار
می‌کند.

\paragraph{بدون نشتی شمارنده.} همهٔ شمارنده‌ها هنگام بارگذاری بسته فقط یک بار تعریف می‌شوند؛ پس در یک سند
می‌توانید هر تعداد تبدیل انجام دهید. آزمونی در مخزن (\tx{leaks.lvt}) همین را بررسی می‌کند.

% ---------------------------------------------------------------------------
\section{محدودیت‌ها و مشکلات شناخته‌شده}\label{sec:limits}

\begin{itemize}
  \item \textbf{ورودی بررسی نمی‌شود.} ماه ۱۳، روز ۴۰ یا رشتهٔ نادرست یا خروجی غلط می‌دهد یا خطای \TeX{}
    (نه پیام خطای خود بسته).
  \item \textbf{ورودی باید خودِ تاریخ باشد.} \cs{pdate}، \cs{gdate}، \cs{pisodate} و \cs{gisodate} ورودی را
    همان‌طور که نوشته شده در خط‌های تیره می‌شکنند و فرمانی را که تاریخ در آن ذخیره شده باز نمی‌کنند. پس
    بنویسید \cs{pdate}\tx{\{2015-03-21\}}، نه \cs{pdate}\cs{mydate}. (\cs{pcalendardatetojulian} ورودی‌اش را
    باز می‌کند و این مشکل را ندارد.)
  \item \textbf{قاعدهٔ کبیسهٔ تقریبی.} قاعدهٔ ۳۳ ساله فقط تقریبی حسابی از تقویم رسمی است که بر پایهٔ رصد
    نجومی تعیین می‌شود. این دو در سال‌های رایج یکی‌اند، اما لزوماً در همهٔ سال‌ها یکی نیستند. برای سال‌های
    دور، نتیجه را با یک جدول معتبر مقایسه کنید.
  \item \textbf{یک اشکال شناخته‌شده.} \cs{pcal@jalalitogregorian}\tx{\{1279-10-11\}} با پیام
    «\lr{Missing number, treated as zero}» متوقف می‌شود، در حالی که تاریخ‌های اطرافش درست‌اند. دلیلش را هنوز
    پیدا نکرده‌ام.
  \item \textbf{حلقه‌های تودرتو.} تبدیل‌ها در داخل خود از \cs{loop} هستهٔ \lr{\LaTeX} استفاده می‌کنند. اگر
    \cs{loop}\ldots\cs{repeat} شما این فرمان‌ها را صدا بزند، بدنهٔ آن را درون یک گروه
    (\cs{begingroup}\ldots\cs{endgroup}) بگذارید، وگرنه حلقه بعد از یک دور متوقف می‌شود. \cs{foreach} این
    مشکل را ندارد.
  \item \textbf{شمارنده‌های مشترک.} هر تبدیل شمارنده‌های نتیجه را از نو می‌نویسد، همان‌طور که در
    بخش~\ref{sec:julian} گفته شد.
  \item \textbf{نام سال} فقط برای ۰ تا ۲۰۹۹ کار می‌کند.
\end{itemize}

% ---------------------------------------------------------------------------
\section{تاریخچهٔ نسخه‌ها}\label{sec:history}

\begin{description}
  \item[\lr{0.5 (2026-11-01)}] رفع نشتی شمارنده (چند \cs{newcount} در هر فراخوانی) و شمارندهٔ تعریف‌نشدهٔ طول
    ماه؛ افزودن این راهنما و \cs{pcalendardate}/\cs{pcalendarversion}.
  \item[\lr{0.4 (2023-11-07)}] آمادهٔ انتشار در \lr{CTAN}.
  \item[\lr{0.3 (2015-08-14)}] نام حیوان سال جلالی، از ۰ تا ۲۰۹۹.
  \item[\lr{0.2 (2015-08-09)}] تبدیل جلالی به میلادی.
  \item[\lr{0.1 (2015-07-25)}] نخستین نسخه.
\end{description}

% ---------------------------------------------------------------------------
\section{مجوز}\label{sec:license}

می‌توانید این اثر را طبق شرایط مجوز عمومی پروژهٔ
\lr{\LaTeX} (\lr{LPPL})، نسخهٔ ۱٫۳ یا هر نسخهٔ جدیدتر به انتخاب خودتان، توزیع یا ویرایش کنید. آخرین نسخهٔ
مجوز در \tx{http://www.latex-project.org/lppl.txt} موجود است. وضعیت نگهداری این اثر در \lr{LPPL}
«\lr{author-maintained}» است، یعنی نویسنده خودش آن را نگهداری می‌کند.

\end{document}
